\documentclass{article}

\usepackage[nonatbib, final]{neurips}
\usepackage[numbers]{natbib}

\makeatletter
\renewcommand{\@noticestring}{\centering}
\makeatother

\usepackage{amsmath,amssymb,amsfonts}
\usepackage{booktabs}
\usepackage{graphicx}
\usepackage{hyperref}
\usepackage{xcolor}
\usepackage{microtype}
\usepackage{multirow}

% Custom macros
\newcommand{\method}{LoRA-Teleport}
\newcommand{\eg}{e.g.,\xspace}
\newcommand{\ie}{i.e.,\xspace}
\newcommand{\etal}{\textit{et al.}\xspace}
\newcommand{\cossim}{\mathrm{cos\_sim}}
\newcommand{\todo}[1]{\textcolor{red}{[TODO: #1]}}

\title{Why Gradient Alignment Teleportation Fails in Sequential Continual Learning:
       A Mechanistic Study}

\author{
  Anonymous Author(s)\\
  \textit{Submitted for blind review}
}

\begin{document}

\maketitle

%% ===================================================================
\begin{abstract}
%% ===================================================================
We present a systematic empirical study of LoRA-based gradient alignment teleportation
for sequential continual learning (CL), adapting the recently proposed COST framework~\citep{zhou2025cost}
from multi-task to sequential settings.
Catastrophic forgetting remains a central challenge in CL, and gradient conflict between
old and new tasks is widely hypothesized as a proximal cause---yet direct manipulation
of gradient alignment during training has not been rigorously tested in the sequential regime.
We implement an online teleportation mechanism that detects gradient conflict at each
training step and applies LoRA perturbations to find loss-equivalent parameter configurations
with better gradient alignment.
Across six hyperparameter configurations on seq-CIFAR-10, all online teleportation variants
\emph{hurt} accuracy relative to a vanilla Experience Replay baseline ($-1.6\%$ to $-31.6\%$).
We identify a structural \emph{zero-order trap}: achieving loss invariance on one batch
requires large parameter perturbations ($\|\Delta\Theta\| = 25$--$30$) that catastrophically
disrupt the training trajectory---better loss invariance produces \emph{worse} accuracy.
A companion experiment further reveals that gradient conflict is \emph{ubiquitous} in 78.9\%
of training steps, providing insufficient discriminating signal to guide teleportation.
These findings constitute a rigorous negative result with a mechanistic explanation
of why LoRA teleportation fails in the sequential regime.
\end{abstract}

%% ===================================================================
\section{Introduction}
\label{sec:intro}
%% ===================================================================

Catastrophic forgetting---the rapid degradation of previously learned knowledge when a
neural network is trained on a new task---remains one of the core obstacles in continual
learning (CL) \citep{kirkpatrick2017ewc, vandeven2022survey}.
A prominent hypothesis attributes forgetting to \emph{gradient conflict}: when the gradient
of the current task opposes the gradient of old tasks, parameter updates that benefit new-task
performance necessarily degrade old-task performance \citep{lopezpaz2017gem, yu2020pcgrad}.

This conflict hypothesis has motivated two main lines of work.
\emph{Constraint-based} methods (\eg GEM~\citep{lopezpaz2017gem}, A-GEM~\citep{chaudhry2019agem})
project current-task gradients to avoid increasing the replay-buffer loss.
\emph{Surgery-based} methods (\eg PCGrad~\citep{yu2020pcgrad}) modify gradient directions
to reduce inter-task conflict.
Both operate at the gradient level---altering how the optimizer steps through parameter space.

A third possibility is \emph{parameter-space teleportation}: instead of modifying the gradient step,
find a loss-equivalent point in parameter space from which the gradient is naturally better aligned.
This idea was formalized by \citet{zhao2022teleport} and extended to multi-task learning (MTL) by
\citet{zhou2025cost} (COST), which uses low-rank adapters (LoRA) to efficiently find such points.
COST demonstrates strong results on simultaneous MTL benchmarks.

\textbf{This paper.}
We ask whether gradient alignment teleportation can reduce catastrophic forgetting in
\emph{sequential} CL, where old tasks are inaccessible and conflict is detected only
through a small replay buffer.
We implement an online variant (\method) that detects gradient conflict during training
and applies LoRA teleportation to reduce it, and conduct a systematic sweep over
hyperparameters on seq-CIFAR-10.

Our main contributions are:
\begin{itemize}
  \item We show that \textbf{all six teleportation configurations hurt accuracy}
        ($-1.6\%$ to $-31.6\%$ vs.\ ER baseline), ruling out hyperparameter explanations
        and pointing to a structural failure.
  \item We identify a \textbf{zero-order trap}: the loss-invariance constraint forces
        parameter perturbations of norm $\|\Delta\Theta\| = 25$--$30$, which disrupt all
        non-target batches and corrupt the training trajectory.
        The configuration with the best loss invariance (avg.\ $\ell_t = 0.46$) produces
        the \emph{worst} accuracy, confirming the structural nature of this failure.
  \item We show that \textbf{gradient conflict is ubiquitous}: 78.9\% of training steps
        exhibit $\cossim(g_\mathrm{old}, g_\mathrm{new}) < 0$,
        consistent with the 50--75\% conflicting neurons found by \citet{liu2026gradconflict}.
        The resulting low variance in per-task mean $\cossim$ (pooled $r = -0.25$, inconsistent
        across seeds) suggests that mini-batch cosine similarity cannot reliably guide teleportation.
\end{itemize}

Together these findings explain \emph{why} LoRA teleportation cannot be directly ported
from MTL to sequential CL, and provide falsifiable, mechanistic accounts of both failure modes.

%% ===================================================================
\section{Background}
\label{sec:background}
%% ===================================================================

\paragraph{Sequential continual learning.}
In class-incremental learning (Class-IL), a model observes tasks $1, \ldots, T$ sequentially.
At test time it must classify across all seen classes with no task identity given.
Experience Replay (ER)~\citep{rolnick2019er} stores a buffer of past samples mixed with
current-task data. We use ER as the base method throughout.

\paragraph{Gradient conflict.}
For parameters $\theta$, conflict between tasks $i$ and $j$ occurs when
$\cossim(g_i, g_j) < 0$ where $g_k = \nabla_\theta \mathcal{L}_k(\theta)$.
GEM~\citep{lopezpaz2017gem} enforces $g^\top g_k \geq 0$ for buffered tasks per step.
PCGrad~\citep{yu2020pcgrad} projects $g_i$ onto the normal plane of $g_j$ when they conflict.
Both modify the gradient direction while leaving $\theta$ unchanged.

\paragraph{Symmetry teleportation.}
Neural networks with ReLU activations have \emph{scaling symmetries}: multiplying pre-activation
weights by $t > 0$ and post-activation weights by $1/t$ preserves network output.
\citet{zhao2022teleport} exploit this to teleport to better-conditioned points on the loss
manifold. COST~\citep{zhou2025cost} extends teleportation to LoRA perturbations in MTL,
minimizing gradient conflict while preserving per-task loss.

\paragraph{LoRA.}
Low-Rank Adaptation~\citep{hu2022lora} parameterizes $\Delta W = BA$ with
$B \in \mathbb{R}^{d \times r}$, $A \in \mathbb{R}^{r \times d}$, $r \ll d$.
After optimization, the LoRA is merged: $W \leftarrow W + BA$, adding no inference overhead.

%% ===================================================================
\section{Method: Online LoRA Teleportation}
\label{sec:method}
%% ===================================================================

At each training step $s$ we draw a current-task batch $(x_\mathrm{new}, y_\mathrm{new})$
and an old-task batch $(x_\mathrm{old}, y_\mathrm{old})$ from a persistent replay iterator.
We compute
\begin{equation}
  \cossim_s = \frac{g_\mathrm{old}^\top g_\mathrm{new}}{\|g_\mathrm{old}\|\|g_\mathrm{new}\| + \varepsilon}
  \label{eq:cossim}
\end{equation}
on the last $L=2$ layers for efficiency.
When $\cossim_s < \tau$ (threshold, default 0.0), we optimize rank-$r$ LoRA perturbations
$\{(A_l, B_l)\}$ on the last $L$ layers via:
\begin{equation}
  \mathcal{L}_\mathrm{teleport} =
    \lambda_t \underbrace{\bigl(|\mathcal{L}(\theta{+}\Delta\Theta, x_\mathrm{old}) - \mathcal{L}(\theta, x_\mathrm{old})|
    + |\mathcal{L}(\theta{+}\Delta\Theta, x_\mathrm{new}) - \mathcal{L}(\theta, x_\mathrm{new})|\bigr)}_{\mathcal{L}_t\text{ (loss invariance)}}
    - \cossim(\theta + \Delta\Theta)
    + \lambda_r \|\Delta\Theta\|_F^2
  \label{eq:teleport_obj}
\end{equation}
where $\Delta\Theta = \sum_l B_l A_l$ (via \texttt{functional\_call}) and
the cosine term uses second-order gradients (\texttt{create\_graph=True}).
After $n$ steps, the LoRA is merged back into $\theta$.
Gradient conflict detection uses the same batch for both detection and teleportation.

\paragraph{Implementation.}
\method\ is built on Mammoth~\citep{buzzega2020mammoth}.
We maintain a \emph{persistent replay iterator} with epoch-based re-shuffling.
A single LoRA state is shared for $\cossim$ computation; an independent copy handles
$\mathcal{L}_t$ to avoid gradient graph interference.
The step counter resets at each task boundary to keep step numbers interpretable.

%% ===================================================================
\section{Experiments: All Teleportation Variants Fail}
\label{sec:h7}
%% ===================================================================

\subsection{Setup}

\textbf{Dataset}: seq-CIFAR-10 (5 tasks, 2 classes/task), Class-IL.
\textbf{Model}: ResNet-18. \textbf{Optimizer}: SGD, lr=0.1, 10 epochs/task, seed=42.
\textbf{Baseline}: ER, buffer=500, achieves 58.53\% Class-IL.
We sweep: \texttt{check\_freq} $\in \{1, 10\}$,
$\lambda_t \in \{1, 10, 50\}$,
$n \in \{5, 10, 20\}$, with rank-2 LoRA on the last 2 layers.

\subsection{Results}

\begin{table}[t]
\centering
\caption{Online \method\ variants on seq-CIFAR-10 (Class-IL accuracy).
         All variants underperform ER.
         \textbf{E6 achieves the lowest avg\_lt yet the worst accuracy} --- the zero-order trap.
         $\Delta$ is the difference from the ER baseline (58.53\%).}
\label{tab:h7}
\setlength{\tabcolsep}{5pt}
\begin{tabular}{llcccccc}
\toprule
Exp & Description & Class-IL & $\Delta$ & Triggers & avg\_lt & $\|\Delta\Theta\|$ \\
\midrule
E0 & ER (baseline)            & \textbf{58.53\%} & --- & 0    & ---  & ---   \\
\midrule
E5 & freq=10, $\lambda_t$=50  & 56.90\%          & $-1.63\%$  & 972  & 1.73 & 9--13 \\
E3 & freq=10, $\lambda_t$=1   & 56.25\%          & $-2.28\%$  & 850  & 5.67 & 4--6  \\
E4 & freq=10, $\lambda_t$=10  & 53.51\%          & $-5.02\%$  & 927  & 1.92 & 9--13 \\
E2 & freq=1,  $\tau$=$-0.3$   & 48.87\%          & $-9.66\%$  & 1834 & 4.80 & ---   \\
E6 & freq=10, $\lambda_t$=10, $n$=20 & 48.22\%   & $-10.31\%$ & 1140 & \textbf{0.46} & 25--30 \\
E1 & freq=1,  $\lambda_t$=1   & 26.93\%          & $-31.60\%$ & 4419 & ---  & ---   \\
\bottomrule
\end{tabular}
\end{table}

Table~\ref{tab:h7} shows all six configurations fall below ER.
The best variant (E5) is still $-1.63\%$ below baseline.
Increasing \texttt{check\_freq} to 1 (check every step) is catastrophic
($-31.6\%$, 12$\times$ training slowdown).

\subsection{The Zero-Order Trap}

\textbf{E6 is the diagnostic case.}
With $\lambda_t = 10$ and $n=20$ steps, $\mathcal{L}_t$ converges to avg\_lt $= 0.46$:
the loss-invariance objective is working as intended.
Yet E6 produces the \emph{worst accuracy} among freq=10 variants ($-10.3\%$).

To understand this paradox, we measure the LoRA delta norm $\|\Delta\Theta\|$ after teleportation.
E6 has $\|\Delta\Theta\| = 25$--$30$, compared to $4$--$6$ for unconstrained variants.
Enforcing $\mathcal{L}_t \approx 0$ on the single target batch requires a large perturbation
that is loss-neutral only at that point.

\paragraph{Why large $\|\Delta\Theta\|$ corrupts training.}
With $3{,}130$ total batches per task, the training trajectory passes through
$3{,}129$ other batches after each teleportation.
A delta $\Delta\Theta$ satisfying $\mathcal{L}(x^*, \theta + \Delta\Theta) = \mathcal{L}(x^*, \theta)$
(zero-order guarantee at $x^*$) has no guarantee for any $x \neq x^*$.
The perturbation at $x^*$ induces an unpredictable, potentially large loss change
on every other batch, effectively injecting noise proportional to $\|\Delta\Theta\|$
into the training trajectory.

\paragraph{The inescapable trade-off.}
Table~\ref{tab:trap} shows that no $\lambda_t$ simultaneously achieves small $\mathcal{L}_t$
\emph{and} small $\|\Delta\Theta\|$ within 5--20 optimization steps with rank-2 LoRA.
This is a structural constraint: in rank-2 LoRA, only $2 \times 2 \times d_L$ degrees of
freedom are available to simultaneously satisfy loss invariance and improve $\cossim$.
The loss landscape around a trained checkpoint is locally non-trivial, and satisfying
two coupled constraints in low rank requires compensating with scale.

\begin{table}[h]
\centering
\caption{The zero-order trap: no regime achieves both small $\mathcal{L}_t$ and small $\|\Delta\Theta\|$.}
\label{tab:trap}
\begin{tabular}{lccc}
\toprule
Regime & avg\_lt & $\|\Delta\Theta\|$ & Net effect \\
\midrule
Unconstrained ($\lambda_t=1$)           & 5.67 & 4--6   & Target batch loss disrupted \\
Moderate ($\lambda_t=50$)               & 1.73 & 9--13  & Partial disruption of trajectory \\
Strongly constrained ($\lambda_t=10$, $n$=20) & \textbf{0.46} & 25--30 & All other batches disrupted \\
\bottomrule
\end{tabular}
\end{table}

%% ===================================================================
\section{Analysis: Is Gradient Conflict the Right Target?}
\label{sec:h8}
%% ===================================================================

Given the failure of teleportation, we examine whether gradient conflict as measured
by mini-batch $\cossim$ is even a reliable signal for forgetting risk.

\subsection{Setup}

We introduce a \emph{detect-only} flag: teleportation is disabled ($\tau = 100$,
never triggered), while $\cossim_s$ is logged at every check step.
We run three seeds (42, 123, 456) on seq-CIFAR-10 ER (10 epochs/task),
\texttt{check\_freq}=10, yielding 313 measurements per eligible task.
We compute per-task forgetting $f_t = \mathrm{peak\_acc}_t - \mathrm{final\_acc}_t$
and Pearson $r$ between mean $\cossim_t$ and $f_t$.

\subsection{Results}

\paragraph{Gradient conflict is ubiquitous.}

\begin{table}[t]
\centering
\caption{Mean $\cossim(g_\mathrm{old}, g_\mathrm{new})$ during seq-CIFAR-10 ER training.
         Parentheses: fraction of steps with $\cossim < 0$.}
\label{tab:h8_cossim}
\begin{tabular}{lccccc}
\toprule
Seed & Task 1 & Task 2 & Task 3 & Task 4 & Mean neg.\ frac.\ \\
\midrule
42  & $-0.173$ (70.6\%) & $-0.252$ (83.1\%) & $-0.214$ (77.0\%) & $-0.133$ (75.1\%) & 76.5\% \\
123 & $-0.179$ (73.2\%) & $-0.216$ (78.6\%) & $-0.304$ (89.1\%) & $-0.155$ (78.3\%) & 79.8\% \\
456 & $-0.234$ (80.8\%) & $-0.221$ (79.6\%) & $-0.268$ (84.7\%) & $-0.131$ (76.4\%) & 80.4\% \\
\midrule
\textbf{Mean} & $\mathbf{-0.195}$ & $\mathbf{-0.230}$ & $\mathbf{-0.262}$ & $\mathbf{-0.140}$ & \textbf{78.9\%} \\
\bottomrule
\end{tabular}
\end{table}

Table~\ref{tab:h8_cossim}: \textbf{78.9\% of all training steps exhibit gradient conflict}.
The mean $\cossim$ varies only from $-0.13$ to $-0.30$ across tasks and seeds (a 0.17-unit range),
leaving very little discriminating signal.

\paragraph{Correlation is weak and inconsistent.}

\begin{table}[h]
\centering
\caption{Pearson $r$ between mean $\cossim$ per task and per-task forgetting.}
\label{tab:h8_corr}
\begin{tabular}{lcc}
\toprule
Seed & $r$(direct: $\bar{c}_t$ vs.\ $f_t$) & $r$(cross: $\bar{c}_t$ vs.\ $f_{t-1}$) \\
\midrule
42  & $-0.43$ & $-0.73$ \\
123 & $-0.01$ & $-0.32$ \\
456 & $-0.54$ & $+0.08$ \\
\midrule
\textbf{Pooled} & $\mathbf{-0.25}$ & --- \\
\bottomrule
\end{tabular}
\end{table}

Table~\ref{tab:h8_corr}: pooled $r = -0.25$, weakly negative but highly inconsistent
(seed 123: $r = -0.01$; seed 456 cross: $r = +0.08$).
Two confounds limit interpretation:
(1) \emph{task-order}: later tasks incur less forgetting and also face slightly less conflict
    (fewer old tasks in the buffer), creating a spurious negative correlation;
(2) \emph{insufficient variance}: the 0.17-unit range of mean $\cossim$ provides too weak
    a signal to separate causal from confounded effects with only 4 tasks per seed.

\paragraph{Implication.}
These results do not disprove the gradient conflict hypothesis; concurrent work
\citep{liu2026gradconflict} provides evidence at the neuron level.
They suggest that \emph{single mini-batch cosine similarity} is not a reliable proxy
for per-task forgetting risk in the sequential setting, and that any teleportation
method gated on this signal will be triggering on noise.

%% ===================================================================
\section{Related Work}
\label{sec:related}
%% ===================================================================

\paragraph{Replay-based CL.}
ER~\citep{rolnick2019er} and DER~\citep{buzzega2020der} mix past samples with current training.
We use ER as our base and augment it with teleportation.

\paragraph{Gradient-based forgetting mitigation.}
GEM~\citep{lopezpaz2017gem} and A-GEM~\citep{chaudhry2019agem} constrain gradient updates
using episodic memory, providing \emph{first-order} per-step guarantees.
Our teleportation seeks a \emph{zero-order} guarantee on a single batch.
The zero-order trap suggests first-order gradient constraints may be more robust:
they modify only the step direction, not the parameter value, and thus do not require
compensating scale to enforce their constraint.

\paragraph{Multi-task gradient surgery.}
PCGrad~\citep{yu2020pcgrad} projects conflicting task gradients in the MTL setting.
MTL's simultaneous task access makes loss computation exact, avoiding the replay-sampling
noise that amplifies the zero-order trap in sequential CL.

\paragraph{Symmetry teleportation.}
\citet{zhao2022teleport} use ReLU symmetries for single-task optimization acceleration.
COST~\citep{zhou2025cost} extends this to MTL via LoRA.
A key distinction from our work: COST evaluates loss on all tasks simultaneously,
whereas sequential CL approximates old-task loss through a small replay buffer.
This approximation error interacts with the zero-order constraint in a way
that COST's MTL setting does not encounter.

\paragraph{Gradient similarity and forgetting.}
\citet{liu2026gradconflict} identify 50--75\% of neurons as conflicting in LLM fine-tuning
and show freezing them reduces forgetting by 59--82\%.
Our aggregate $\cossim$ finding (78.9\% conflicting steps) is consistent with their
neuron-level statistics, but our step-level correlation analysis suggests that coarser
mini-batch $\cossim$ does not have the resolution needed for step-by-step teleportation decisions.

\paragraph{LoRA in CL.}
Several concurrent methods use LoRA for parameter isolation in CL
\cite{PLACEHOLDER_treelora, PLACEHOLDER_keeplora}, but do not use teleportation.
Their positive results suggest LoRA is a useful parameterization for CL,
but the specific objective of gradient alignment via teleportation is what fails in our work.

%% ===================================================================
\section{Limitations}
\label{sec:limitations}
%% ===================================================================

\paragraph{Benchmark scope.}
All experiments use seq-CIFAR-10 (5 tasks, ResNet-18).
Generalization to longer task sequences, transformer architectures, or LLM fine-tuning is not confirmed.

\paragraph{Low-rank budget.}
Rank-2 LoRA with 5--20 steps may be insufficient to find the desired $\Delta\Theta$.
Higher ranks or more steps might reduce $\|\Delta\Theta\|$---at substantially higher compute cost
(second-order gradients scale quadratically with rank).

\paragraph{H8 statistical power.}
With 4 eligible tasks per seed, the correlation analysis has 12 pooled data points.
Stronger tests would use longer sequences or per-step analysis.

\paragraph{Alternative conflict signals.}
Cumulative cosine similarity over a trajectory window, or per-neuron conflict
as in \citet{liu2026gradconflict}, might provide more reliable teleportation triggers.

%% ===================================================================
\section{Conclusion}
\label{sec:conclusion}
%% ===================================================================

We present a rigorous negative result: LoRA-based gradient alignment teleportation fails
in sequential continual learning.
Two structural barriers account for the failure.

\textbf{The zero-order trap} (\S\ref{sec:h7}): Loss invariance on a single target batch forces
$\|\Delta\Theta\| = 25$--$30$, disrupting the training trajectory through all non-target batches.
The configuration with the best loss invariance produces the worst accuracy.
This creates an inescapable dilemma: either enforce loss invariance and corrupt the trajectory,
or relax it and allow the optimizer to sacrifice it for $\cossim$ gain.

\textbf{Ubiquitous conflict} (\S\ref{sec:h8}): Gradient conflict is present in 78.9\% of
training steps. The low variance in mean $\cossim$ across tasks (0.17 units) and weak,
inconsistent correlation with forgetting ($r = -0.25$, high seed variance) suggest that
mini-batch $\cossim$ is not a reliable signal for driving teleportation in sequential CL.

Our findings do not contradict the broader gradient conflict hypothesis---\citet{liu2026gradconflict}
present compelling neuron-level evidence, and first-order methods like GEM/A-GEM have demonstrated
value. Rather, they identify \emph{parameter-space teleportation} as an ineffective vehicle
for conflict reduction in the sequential regime. Future work might explore first-order gradient
projection adapted for sequential CL, or neuron-level selective freezing as in \citet{liu2026gradconflict},
as more principled alternatives.

\bibliographystyle{plainnat}
\bibliography{bibliography}

\end{document}
